\section{The increment of capital supply}\label{sec:inc}

Section~\ref{sec:mfg} ends by closing a route and leaving the original target: uniqueness through
single crossing, which the programme has pursued through Light's Theorem 1 and cannot get above
$\mu\approx3.4$. This section changes what has to be proved. Theorem 1 need not hold. It need only
fail on little mass.

\subsection{What single crossing actually asks}

Work in wage units, where homotheticity makes capital supply per unit of wage a function of the rate
alone and clearing reads $S(r)=D(r)$ with $D(r)=\alpha/((1-\alpha)(r+\delta))$. Excess supply is
increasing --- which is all single crossing needs --- as soon as
\begin{equation}\label{eq:inc}
  S(r_2)-S(r_1)\;>\;D(r_2)-D(r_1),\qquad r_1<r_2,
  \tag{I}
\end{equation}
that is, as soon as supply does not fall \emph{faster} than demand. Demand is steep at the
calibration, and that is the whole opportunity: over the proved bracket $r\in[2\%,8\%]$ supply's
worst secant slope is $+33$ at $\gamma=5$ against $-54$ for demand, and it never bends backwards at
any $\gamma$ up to $12$. A proof of \eqref{eq:inc} may be lossy by more than the true increment and
still close.

The two-sided bracket on the \emph{level} of capital that the existence theorem supplies cannot be
used for this. As $r_2-r_1\to0$ the bracket-only condition $F(r_2)-C(r_1)>D(r_2)-D(r_1)$ degenerates
into $F\ge C$, a bracket of zero width; measured, its margin runs from $-352$ to $-1737$. The object
to bound is the increment, not the level twice.

The two-sided bracket on the \emph{level} of capital that the existence theorem supplies
cannot be used for this, for a reason that recurs below and is set out in
Appendix~\ref{app:routes}: a bracket as wide as the quantity being bounded contributes a
spurious spread larger than the real one.

\subsection{Coupling the two economies}

Run both economies on a common earnings path from the common start $a_0=0$, and let
$x_j=a_j(r_2)-a_j(r_1)$. Split the step at the \emph{lower}-rate state:
\[
  x_{j+1}=\underbrace{\bigl[g^2_j(a^2_j)-g^2_j(a^1_j)\bigr]}_{\text{same policy, different assets}}
   +\underbrace{\bigl[g^2_j(a^1_j)-g^1_j(a^1_j)\bigr]}_{\text{same state, different rate}} .
\]
The second bracket is the failure of Theorem 1 at that state, which
(\lean{stagePolicy\_le\_iff\_consumption\_le}) is exactly $c^2-c^1\le(r_2-r_1)a$ read with a slack;
call its shortfall $e_j$. The first is at least $-L_j\,(x_j)^-$ when saving is nondecreasing in
assets and Lipschitz with constant $L_j$. So
\begin{equation}\label{eq:rec}
  x_{j+1}\;\ge\;-L_j\,(x_j)^-\;-\;e_j .
  \tag{$\ast$}
\end{equation}
Nothing in \eqref{eq:rec} asks $x_j\ge0$: the pointwise ordering of policies is not assumed, and its
failure is priced by $e_j$ rather than excluded.

\begin{theorem}[The gap envelope]\label{thm:envelope}
If $x_0\ge0$, $L_j\ge0$, $e_j\ge0$ and \eqref{eq:rec} holds at every age, then $x_j\ge-f_j$ where
$f_0=0$ and $f_{j+1}=L_jf_j+e_j$; hence $\sum_jx_j\ge-\sum_jf_j$. With demand falling by more than
the average of $f$, \eqref{eq:inc} follows.
\formal{OLG.gapAccum}, \lean{OLG.negPart\_le\_gapAccum}, \lean{OLG.neg\_gapAccum\_le},
\lean{OLG.sum\_neg\_gapAccum\_le}, \lean{OLG.lt\_of\_increment\_bound}
\end{theorem}

\subsection{The mass-weighted form}

The negative part is the point of the formulation. From \eqref{eq:rec},
\[
  (x_{j+1})^-\;\le\;L_j\,(x_j)^-+e_j ,
\]
whose right-hand side is \emph{linear} in the two quantities. Expectations therefore pass straight
through it, and averaging over the cohort replaces $e_j$ by the \textbf{mass-weighted} failure of
Theorem 1, $\bar e_j=\mathbb E_{1}\bigl[(g^1_j-g^2_j)^+\bigr]$, in place of its worst case over
states. Since the violation is confined to a sliver of the distribution, that is the difference
between a bound that fails and one that passes by two orders of magnitude.

Measured at the calibration, with $L_j=(1-\kappa_k)(1+r_2)$ from the sandwich and the accumulation
$\Gamma=\frac1J\sum_jf_j/\varepsilon$ around $21$ to $24$:

\begin{center}
\begin{tabular}{lrrrrr}
\toprule
& \multicolumn{2}{c}{$r\in[4\%,6\%]$, budget $0.670$}
& \multicolumn{2}{c}{$r\in[2\%,8\%]$, budget $2.109$}\\
\cmidrule(lr){2-3}\cmidrule(lr){4-5}
$\gamma$ & $5$ & $8$ & $5$ & $8$\\
\midrule
worst violation $\varepsilon$ & $0.0460$ & $0.0701$ & $0.1313$ & $0.2053$\\
mass carrying it & $0.0091$ & $0.0666$ & $0.0084$ & $0.0636$\\
\midrule
bound, uniform $\varepsilon$ & $0.976$ $\times$ & $1.464$ $\times$ & $3.125$ $\times$ & $4.707$ $\times$\\
bound, age-dependent $\varepsilon_j$ & $0.760$ $\times$ & $1.218$ $\times$ & $2.498$ $\times$ & $3.953$ $\times$\\
bound, mass-weighted $\bar e_j$ & $\mathbf{0.0046}$ $\checkmark$ & $\mathbf{0.0310}$ $\checkmark$ & $\mathbf{0.0138}$ $\checkmark$ & $\mathbf{0.0958}$ $\checkmark$\\
\midrule
true $S(r_2)-S(r_1)$ & $+0.692$ & $+0.457$ & $+2.060$ & $+1.368$\\
\bottomrule
\end{tabular}
\end{center}

At $\gamma\le3$ every column is zero: in wage units Theorem 1 holds exactly, so $\varepsilon=0$ and
the bound is trivial. The interest is at $\gamma=5$ and $8$, where Theorem 1 fails and the
mass-weighted bound passes with factors of $145$ and $22$ to spare --- and passes on the wide
bracket too, where the rates differ by six percentage points. Charging the worst violation at every
age fails; charging it where it occurs does not. Age-dependence alone does not rescue the worst-case
version, because the violation is spread over $53$ of $60$ ages while being tiny at nearly all of
them.

\subsection{The Lipschitz constant}

The engine is proved, and so is the Lipschitz constant it runs on.

\subsubsection*{The propensity floor}

\begin{theorem}[The propensity to consume is at least $\kappa_k$]\label{thm:mpcfloor}
At every stage, every income state and every pair of asset levels in the region,
\[
  \kappa_k\,(1+r)(a_2-a_1)\;\le\;c_k(z,a_2)-c_k(z,a_1),\qquad a_1\le a_2 ,
\]
so saving is Lipschitz in assets with constant $(1-\kappa_k)(1+r)$, which is the $L_j$ the increment
argument uses.
\formal{IncomeFluctuation.stageMPC\_mul\_le\_stageConsumption\_sub}
\end{theorem}

The terminal condition is an identity rather than an assumption: in the last period consumption is
cash on hand, so the propensity is exactly $1=\kappa_0$. The step is the two-point Euler comparison.
With $b_i$ the saving at $a_i$ and $t=\kappa_kR(b_2-b_1)\ge0$, the induction hypothesis raises next
period's consumption by at least $t$ in \emph{every} state, so the Euler equation above the floor at
$a_2$ gives
\[
  c_2^{-\gamma}\;\le\;\beta R\sum_{z'}\pi\,\bigl(C'(z')+t\bigr)^{-\gamma}
   \;=\;\beta R\,M(C'+t)^{-\gamma}\;\le\;\beta R\,\bigl(M(C')+t\bigr)^{-\gamma},
\]
with $M$ the certainty equivalent, the power mean of exponent $-\gamma$. The Euler equation under
the cap at $a_1$ gives $M(C')\ge\Thorn c_1$, and $\beta R=\Thorn^\gamma$ turns the chain into
$c_2\ge c_1+t/\Thorn$, which is the claim at stage $k+1$ once denominators are cleared.

The last inequality in the display is \textbf{superadditivity of the certainty equivalent},
$M(x+t)\ge M(x)+t$, which is concavity and homogeneity of degree one at the midpoint
(\lean{powerMean\_add\_const\_le}, already proved for the constraint-incidence bound of
Section~\ref{sec:thm1}). It is worth pausing on: this is the direction a negative power mean
obliges, and it is the reverse of what several arguments in this programme wanted and could not
have. The recurring obstruction, read the other way round, is what makes the floor go through.

Measured, the floor holds at every stage and every reached state with equality only at the last
stage, and the constant is nearly tight: the true accumulation is $20.7$ against the sandwich's
$21.2$.

\subsubsection*{The Lipschitz bound it gives}

\begin{theorem}[Accumulated assets are Lipschitz in assets]\label{thm:cohortlip}
For $a_1\le a_2$ in the region,
\[
  A_k(a_2,z)-A_k(a_1,z)\;\le\;G_k\,(a_2-a_1),
\]
with $G_k$ the slope \lean{cohortSlope} of the affine ceiling: $G_0=1$,
$G_{k+1}=1+G_k(1-\kappa_{k+1})(1+r)$.
\formal{IncomeFluctuation.cohortAssets\_sub\_le\_cohortSlope}
\end{theorem}

The induction is the one already written for the affine floor, with one substitution: the propensity
floor turns a step of assets into a step of saving no larger than $(1-\kappa_{k+1})(1+r)$ times it,
and the expectation over next period's earnings passes the constant through unchanged. The same $G_k$
that bounds the \emph{level} of accumulated assets in Section~\ref{sec:eqm} therefore bounds their
\emph{differences}, which is what the increment argument accumulates --- and it is the $L_j$ whose
accumulation was measured at $21$ above.

\subsection{The mass-weighted failure}

What is left is $\bar e_j$, an expectation against a distribution nobody controls. It looks like the
hard part, and it would be if the failure of Theorem 1 were spread over assets. It is not.

Measured, the failure of Theorem 1 is not spread over assets but confined to the
high-earnings states --- at $\gamma=5$ the top two of seven, carrying invariant mass $0.095$.
That is the coordinate in which no estimate is needed.

\subsubsection*{Why that coordinate is free}

Earnings are exogenous and the cohort is born with its state drawn from the invariant law, so the
earnings marginal is $\nu$ at \emph{every} age, exactly. Nothing has to be estimated in that
coordinate.

\begin{theorem}[The earnings marginal of a cohort is $\nu$]\label{thm:marginal}
If $\nu$ is invariant for $\pi$, the $\nu$-weighted expectation of a function of the earnings state
is unchanged by any number of steps of the cohort recursion. Hence
$\bar e_j\le\sum_z\nu_z\,V_j(z)$ with $V_j(z)$ the supremum over assets of the failure at earnings
state $z$: exact in earnings, worst-case only in assets.
\formal{IncomeFluctuation.incomeStep}, \lean{IncomeFluctuation.stageStep\_of\_incomeFn},
\lean{IncomeFluctuation.sum\_mul\_incomeStep}, \lean{IncomeFluctuation.sum\_mul\_incomeStep\_iterate},
\lean{IncomeFluctuation.sum\_mul\_le\_of\_le}
\end{theorem}

Being exact in the coordinate where the failure lives is what makes the bound work:

\begin{center}
\begin{tabular}{lrrrr}
\toprule
& \multicolumn{2}{c}{$r\in[4\%,6\%]$, budget $0.670$}
& \multicolumn{2}{c}{$r\in[2\%,8\%]$, budget $2.109$}\\
\cmidrule(lr){2-3}\cmidrule(lr){4-5}
$\gamma$ & $5$ & $8$ & $5$ & $8$\\
\midrule
worst-case in both coordinates & $0.976$ $\times$ & $1.464$ $\times$ & $3.125$ $\times$ & $4.707$ $\times$\\
Markov in assets, best threshold & $0.976$ $\times$ & $1.464$ $\times$ & $3.125$ $\times$ & $4.707$ $\times$\\
$\nu$-exact, worst-case in assets & $\mathbf{0.032}$ $\checkmark$ & $\mathbf{0.130}$ $\checkmark$ & $\mathbf{0.103}$ $\checkmark$ & $\mathbf{0.422}$ $\checkmark$\\
fully mass-weighted (the truth) & $0.0046$ & $0.0310$ & $0.0138$ & $0.0958$\\
\bottomrule
\end{tabular}
\end{center}

The $\nu$-exact row passes with factors of $21$, $5.1$, $20$ and $5.0$, and it is within a factor of
seven of the fully mass-weighted truth --- the whole of the remaining looseness being the worst case
over assets inside each earnings state.

\subsection{The two-rate difference}

One quantity remains: $V_j(z)$, the largest failure of Theorem 1 over assets at a single earnings
state. There is no distribution in it, and the accumulation leaves an allowance of
$0.670/21\approx0.032$ per age for $\sum_z\nu_zV_j(z)$. Two routes to it suggest themselves and both
fail --- a quantitative version of the Euler step, which amplifies, and a direct bound from the
sandwich, which is an order of magnitude too loose (Appendix~\ref{app:routes}). What works is a
different chaining of the same two Euler relations.

\subsubsection*{Chaining through the certainty equivalent}

That amplification is not intrinsic to comparing the two rates; it is an artefact of how the
expectation was bounded. The minimum-based step \eqref{eq:qstep} of Appendix~\ref{app:routes}
replaces $\sum\pi\,C^{-\gamma}$ by its smallest term,
which charges every stage's error against next period's consumption in the \emph{worst} earnings
state. Chaining the same two Euler relations through the certainty equivalent instead leaves no
minimum anywhere:

\begin{proposition}[Certainty-equivalent chaining]\label{prop:cechain}
From $c_2^{-\gamma}\ge K_2M_2^{-\gamma}$ and $c_1^{-\gamma}\le K_1M_1^{-\gamma}$,
\[
  \frac{c_2}{c_1}\;\le\;\Bigl(\frac{K_1}{K_2}\Bigr)^{1/\gamma}\frac{M_2}{M_1},
\]
with $M_i$ the certainty equivalent of next period's consumption.
\formal{CertaintyEquivalent.mul\_le\_mul\_of\_euler\_pair}
\end{proposition}

Propagating a constant shift through $M$ then costs the power mean's gradient sum
$A=\sum_{z'}\pi\,(C/M)^{-\gamma-1}$, which is $1$ exactly when next period's consumption is riskless.
So the per-stage multiplier becomes a measure of the dispersion of next period's consumption under the
risk-tilted measure, rather than the ratio of today's consumption to its worst realisation:

\begin{center}
\begin{tabular}{lrrrrr}
\toprule
$\gamma$ & $1$ & $2$ & $3$ & $5$ & $8$\\
\midrule
multiplier, minimum-based & $5.75$ & $5.47$ & $5.20$ & $4.72$ & $4.20$\\
multiplier $A$, certainty-equivalent & $\mathbf{1.041}$ & $\mathbf{1.063}$ & $\mathbf{1.086}$ & $\mathbf{1.132}$ & $\mathbf{1.189}$\\
$qA$, with $q=(R_1/R_2)^{1/\gamma}$ & $1.022$ & $1.053$ & $1.079$ & $1.127$ & $1.186$\\
$(qA)^{59}$ & $3.6$ & $21$ & $88$ & $1180$ & $2.2\times10^{4}$\\
\bottomrule
\end{tabular}
\end{center}

A fivefold reduction per stage, and at $\gamma=1$ the sixty-stage accumulation falls from
$5.75^{59}$, which is beyond any use, to $3.6$. It is still not a contraction: $qA>1$ throughout, so
the recursion grows, and the per-stage criterion it yields,
$b_1\le(\Thorn/qA)\bigl(a+c_1/(\gamma R)\bigr)$ --- which is the \lean{StageSlack} of
Section~\ref{sec:thm1} with the risk penalty $A$ made explicit --- still fails by fourteen per cent at
zero wealth in the top earnings state at $\gamma=1$.

\subsubsection*{The propensity floor damps the step}

The factor of $1.02$ closes once the propensity floor is put into the step, where it belongs. Under
the failure hypothesis the household at the lower rate saves MORE, by exactly the failure $D$, and
saving more raises next period's consumption by at least $\kappa_kR_1D$
(Theorem~\ref{thm:mpcfloor}). Carrying that through,
\[
  D\Bigl[1+\frac{qA\kappa_kR_1}{\Thorn_1}\Bigr]\;\le\;\frac{qA}{\Thorn_1}\varepsilon_k
   +(q-1)c_1+\frac{qA}{\Thorn_1}\Delta r\,b_2-\Delta r\,a ,
\]
so the error carried from one stage to the next is multiplied not by $qA/\Thorn_1$ but by
\[
  m_k\;=\;\frac{qA/\Thorn_1}{1+qA\kappa_kR_1/\Thorn_1}\;=\;\frac{qA}{\Thorn_1+qA\kappa_kR_1},
\]
at most one exactly when $qA\le\Thorn_1+qA\kappa_kR_1$
(\lean{CertaintyEquivalent.damped\_le\_one}). The damping is the propensity floor feeding back
against the failure.

\begin{center}
\begin{tabular}{lrrrrrr}
\toprule
$\gamma$ & $1$ & $1.5$ & $2$ & $2.5$ & $3$ & $5$\\
\midrule
$\max_k m_k$ & $\mathbf{0.978}$ & $\mathbf{0.992}$ & $1.002$ & $1.012$ & $1.021$ & $1.035$\\
$\prod_k m_k$ & $0.027$ & $0.062$ & $0.112$ & $0.196$ & $0.337$ & $0.785$\\
\bottomrule
\end{tabular}
\end{center}

At $\gamma\le1.5$ the recursion contracts stage by stage. And the product over the sixty stages is
below one at \emph{every} risk aversion tested, because $\kappa_k$ approaches one late in life, where
the damping is strongest --- so an error carried from the terminal condition, which is exactly zero, is
damped over a life even at $\gamma=5$. This is the first non-amplifying propagation of the two-rate
difference in this programme.

\subsubsection*{The multiplier is a variance}

The multiplier $A$ has to be bounded from primitives, and what does that is the exact constraint
$M$ carries rather than any bound on the spread of $C$ (Appendix~\ref{app:routes} collects the
attempts that use the spread). Writing $u=(C/M)^{-\gamma}$ for the risk
tilt, the definition of the power mean says $\mathbb E_w[u]=1$ \emph{exactly}, and the gradient sum
is $A=\mathbb E_w[u^{1+1/\gamma}]$. At $\gamma=1$ the exponent is two and

\begin{proposition}[The gradient sum is one plus the variance]\label{prop:gradvar}
With $\mathbb E_w[u]=1$, $\;\mathbb E_w[u^2]=1+\operatorname{Var}_w(u)$.
\formal{sum\_mul\_sq\_eq\_one\_add\_variance}
\end{proposition}

so $A=1+\operatorname{Var}_w(u)$, with no inequality anywhere. Above $\gamma=1$ the exponent falls
below two and Lyapunov gives $A\le(1+\operatorname{Var}_w(u))^{(\gamma+1)/(2\gamma)}$.

\begin{center}
\begin{tabular}{lrrrrr}
\toprule
$\gamma$ & $1$ & $1.5$ & $2$ & $2.5$ & $3$\\
\midrule
allowance on $A$ & $1.066$ & $1.060$ & $1.056$ & $1.054$ & $1.053$\\
$A$ true & $1.041$ & $1.050$ & $1.058$ & $1.067$ & $1.076$\\
$\operatorname{Var}_w(u)$ & $0.041$ & $0.092$ & $0.165$ & $0.270$ & $0.416$\\
bound from the variance & $\mathbf{1.041}$ & $1.076$ & $1.122$ & $1.182$ & $1.261$\\
slack on the variance & $\mathbf{1.61\times}$ & $0.78\times$ & $0.46\times$ & $0.29\times$ & $0.19\times$\\
\bottomrule
\end{tabular}
\end{center}

At $\gamma=1$ the bound is exact and leaves $1.61$ times of room on the variance at the newborn stage
--- $2.35$ times at stage twenty --- so the whole chain closes there on one remaining primitive input:
$\operatorname{Var}_w\bigl((C/M)^{-1}\bigr)\le0.066$, against a measured $0.041$. Above $\gamma=1$ it
does not: Lyapunov loses accuracy as the exponent falls below two, and the variance itself grows
tenfold by $\gamma=3$, so the bound exceeds the allowance from $\gamma=1.5$ on.

That is a complete route at logarithmic utility, and a sharp statement of what it needs. It would give
Theorem 1 at \emph{every} state, not only the binding one, hence monotone capital supply and
uniqueness at $\mu=1$ --- which is more than Proposition~\ref{prop:twosand} delivers.

\subsubsection*{The two-regime bracket}

So the variance of the risk tilt is what must be bounded, and bounding it needs an admissible set
for $C$ tight enough that its worst member satisfies the requirement. Read as a bound, the constrained
regime's identity supplies a \emph{ceiling}: consumption never exceeds cash on hand, so at every
state
\[
  \underbrace{\kappa_k\bigl(m+H^{\mathrm{risk}}_k(z')\bigr)}_{L(z')}\;\le\;C(z')\;\le\;
  \underbrace{\min\bigl\{\kappa_k(m+H_k(z')),\;m\bigr\}}_{U(z')},\qquad m=y(z')+Rb,
\]
both sides primitive, together with the slope constraints $\kappa_k\Delta y\le\Delta C\le\Delta y$ that
link the states. Maximising $\operatorname{Var}_w(M/C)$ over that admissible set changes the picture
completely. The best feasible profile found gives $0.0406$ at the newborn stage --- essentially the
true variance of $0.0404$ --- against an allowance of $0.0661$, and a search over the set turns up
nothing above it. The slope constraints are what do the work: dropping them and maximising over the
box alone gives $0.1537$.

It is not a certificate, and the gap is worth being exact about. What is established rigorously is an
\emph{interval}: the true consumption profile is feasible, so the maximum is at least $0.0411$; the
polytope sits inside the box, so the maximum is at most $0.1537$. The allowance of $0.0661$ lies
inside that interval. The search says the maximum is at the lower end and the chain closes; the
search is not exhaustive, and at two of the three stages tested it returns a value below the true
variance, which shows directly that it can miss.

Closing this needs an upper bound on the maximum that uses the weights. The weight-free one is
available and useless: since $\mathbb E_w[u]=1$, $\operatorname{Var}_w(u)\le\max u-1\le
\mathrm{HM}(U)/\min L-1$, which is $0.97$ at the binding stage against an allowance of $0.066$.

\subsubsection*{The coefficient of variation}

The certificate wants an upper bound on the maximum, and the coordinates matter for getting one. Put $v_i=1/C_i$. Put $v_i=1/C_i$. Then $M=1/\sum_jw_jv_j$ and $u_i=Mv_i$, so
\begin{equation}\label{eq:cv}
  \operatorname{Var}_w(u)\;=\;\frac{\sum_iw_iv_i^{2}}{\bigl(\sum_iw_iv_i\bigr)^{2}}-1 ,
  \tag{CV}
\end{equation}
the squared coefficient of variation of $v$ under $w$ --- exact algebra, with the mean eliminated
rather than bounded. The box in $C$ is a box in $v$, the slope constraints are still linear in $C$,
and \eqref{eq:cv} is a ratio of quadratic forms, scale-invariant, of the kind whose maximum over a
polytope is a standard finite computation.

\subsubsection*{Why enumeration is legitimate}

In $v$-coordinates the objective is coordinatewise quasi-convex, and that is what makes enumeration
legitimate. Holding all but one coordinate fixed,
\[
  \frac{\partial}{\partial v_j}\frac{\sum_iw_iv_i^{2}}{(\sum_iw_iv_i)^{2}}
  \;\propto\;v_jB-A,\qquad A=\sum_{i\ne j}w_iv_i^{2},\quad B=\sum_{i\ne j}w_iv_i,
\]
so the objective is unimodal with a \emph{minimum} in each coordinate: its maximum over any interval
is at an endpoint. At a maximiser every coordinate therefore sits at an endpoint of its conditional
feasible range.

\subsubsection*{Concavity is the missing constraint}

The bracket and the slope constraints are not by themselves enough: enumerating the admissible set
turns up a profile with variance $0.0874$, above the allowance of $0.0661$
(Appendix~\ref{app:routes}). The profiles that do the damage zig-zag across earnings states, and they
are not consumption functions: consumption is \emph{concave} in cash on hand, by Carroll--Kimball, which this development
already proves (\lean{concaveOn\_stageConsumption}). Across earnings states at fixed assets, cash on
hand moves with $y$, so concavity says the slopes $\Delta C/\Delta y$ fall with $j$ --- five more
linear constraints that the bracket was silently discarding.

With them the same two-directional enumeration gives $0.0415$ at every stage tested, against
allowances of $0.0661$, $0.0725$ and $0.0969$: a margin of $1.59$ to $2.34$ times.

\begin{center}
\begin{tabular}{lrrr}
\toprule
stage & $59$ & $40$ & $20$\\
\midrule
allowance & $0.0661$ & $0.0725$ & $0.0969$\\
enumerated, bracket $+$ slopes & $0.0874$ $\times$ & $0.0702$ & $0.0415$\\
enumerated, $+$ concavity & $\mathbf{0.0415}$ & $\mathbf{0.0415}$ & $\mathbf{0.0415}$\\
\bottomrule
\end{tabular}
\end{center}

So the ingredient the chain was missing is a theorem it already has.

\subsubsection*{The certificate in exact arithmetic}

A first run of this certificate in floating point, over a grid of asset values, leaves two gaps:
the arithmetic, and the sampling (Appendix~\ref{app:routes}). Both close. The calibration is rebuilt
so that every constant is rational: the
earnings chain is Rouwenhorst rather than Tauchen, whose transition probabilities are rational
functions of $\rho=9/10$ and whose invariant law is binomial, and the earnings levels are rational,
normalised to mean exactly one. At $\gamma=1$ everything downstream stays rational ---
$\Thorn=\beta R_1=624/625$, $q=R_1/R_2=52/53$, the propensities $\kappa_k$ by their recursion,
arithmetic human wealth by averages and risk-adjusted human wealth by harmonic means --- so the
allowance is an exact rational number,
\[
  \frac{\Thorn}{q\,(1-\kappa_{J-1}R_1)}-1\;=\;0.066143400439\ldots
\]
The computation is carried in integers over $10^{12}$ with every operation rounded \emph{outward}, so
each quantity is a true enclosure; and the asset argument is carried as intervals covering its whole
range rather than sampled, which at the newborn stage is $b\in[0,3]$, since a newborn holds nothing and
cannot save more than $y_{\max}$.

The branch and bound certifies every node: seven earnings states, the asset range split into four
intervals, at most $3385$ boxes at the worst of them. Nothing is sampled and nothing is rounded
inward.

It survives a much larger asset range than the newborn needs. Rerun over $b\in[0,12]$ in eight
intervals --- four times the range any newborn can reach, and beyond what the later stages of the life
cycle require --- every node is still certified, with at most $16673$ boxes. So the asset argument is
not where the remaining scope question lies.

\subsubsection*{Every stage}

Running the certificate at each stage in turn --- each with its own $\kappa_k$, its own arithmetic and
risk-adjusted human wealth, its own allowance, and the asset argument as intervals over $[0,12]$ in
four pieces --- certifies all fifty-nine of them. The allowance falls monotonically from $1.168$ at
$k=1$ to $0.066143$ at the newborn stage, and every node is discarded at every one.

The cost is where the difficulty lives, and it does not plateau. The early stages are nearly free ---
stage $1$ needs a single box, stage $5$ needs $539$ --- and the count then climbs from about $3600$
around stage $40$ to $14803$ at stage $44$ and $32733$ at stage $59$, a tenfold growth over the last
twenty stages as the allowance tightens towards the variance. The whole run takes about ten hours.

\subsubsection*{What is settled}

At the rational calibration, at logarithmic utility, over $r\in[4\%,6\%]$: the variance is certified
below what the damping allows, in exact arithmetic, at every stage and over the whole asset range.
With the eight formalised lemmas of this section and Section~\ref{sec:thm1}, Theorem 1 holds at every
state and every age; capital supply rises with the rate; the stationary equilibrium is unique.

What that is, precisely, is a theorem with a computational hypothesis discharged by certificate rather
than by proof term --- the same standing as the existence bracket of Section~\ref{sec:eqm}, and a
weaker standing than the lemmas it rests on. The calibration is one calibration: $\beta=24/25$,
$\rho=9/10$, seven earnings states, sixty periods, $r\in[4\%,6\%]$, logarithmic utility. Nothing here
extends to $\gamma>1$, where Section~\ref{sec:thm1} shows the certainty-equivalent side of the
argument fails from $\gamma\approx1.5$ on. The log case is closed; the rest is not. The certificate is
\texttt{numerics/rational/cert.py}, with the calibration in \texttt{inputs.py}. Numerics: \texttt{cestep.m}, \texttt{damp.m},
\texttt{nearfar.m}, \texttt{lyap.m}, \texttt{varbnd.m}, \texttt{varmpc.m}, \texttt{varmax.m},
\texttt{slopestate.m}, \texttt{tworeg.m}, \texttt{polymax.m}, \texttt{bandb.m}, \texttt{vertcert.m},
\texttt{vertfull.m}, \texttt{certify.m}.

What is proved stands on its own: single crossing needs only the increment \eqref{eq:inc}; the
increment is controlled by the envelope of Theorem~\ref{thm:envelope}; its Lipschitz constant is
Theorem~\ref{thm:mpcfloor} and Theorem~\ref{thm:cohortlip}; and the mass weighting is exact in the
coordinate where the failure lives, by Theorem~\ref{thm:marginal}. Numerics: \texttt{incbound.m},
\texttt{ebar.m}, \texttt{vjbound.m}, \texttt{vjdiag.m}, \texttt{mpclow.m}, \texttt{slopecheck.m}.
